\section{Block-diagonal boundary closing}\label{sec:block_diagonal_boundary_closing}

\begin{lemma}[Normalized BNT blocks give the large-length block intersection]
    \label{lem:bnt_direct_sum_unital_block_intersection_ge}
    \lean{MPSTensor.pgvwc07_directSum_restriction_intersection_of_wordTupleSpanTop}
    \lean{MPSTensor.pgvwc07_iSup_restriction_intersection_of_ge_of_bnt_directSum_unital_c1}
    \lean{MPSTensor.pgvwc07_directSum_restriction_intersection_of_ge_of_bnt_directSum_unital_c1}
    \leanok
    \uses{lem:unital_block_separating_product_span,
        lem:product_span_block_image_direct_sum,
        lem:block_restriction_intersection_subspace}
    Let \(A^1,\ldots,A^r\) be separated BNT blocks satisfying the irreducibility,
    left-canonicality, and normalized self-overlap hypotheses used in the
    direct-sum comparison. Suppose \(L_0>0\), and for every block \(j\) the map
    \[
        \Gamma_{L_0}^{A^j}:M_{D_j}(\mathbb C)\to
        (\mathbb C^d)^{\otimes L_0}
    \]
    is injective. Assume also the normalization
    \[
        \sum_a A^j_aA^{j\,\dagger}_a=I .
    \]
    Set
    \[
        N=(L_0+1)+(r-1)3(L_0+1).
    \]
    For every \(n\ge N\), set
    \[
        S_n:=\bigvee_jG_{n+1}(A^j),
    \]
    and the two consecutive sums are internal direct sums:
    \[
        \bigvee_jG_{n+1}(A^j)
        =
        \bigoplus_jG_{n+1}(A^j),
        \qquad
        \bigvee_jG_{n+2}(A^j)
        =
        \bigoplus_jG_{n+2}(A^j).
    \]
    With this notation,
    \[
        \left\{\psi:\forall b,\ \psi(-,b)\in S_n\right\}
        \cap
        \left\{\psi:\forall a,\ \psi(a,-)\in S_n\right\}
        =
        \bigvee_jG_{n+2}(A^j).
    \]
\end{lemma}

\begin{proof}
    \leanok
    Lemma~\ref{lem:unital_block_separating_product_span} gives
    \[
        \operatorname{span}\{(A^1_w,\ldots,A^r_w):|w|=n\}
        =
        \prod_jM_{D_j}(\mathbb C)
    \]
    for every \(n\ge N\). Substituting this product-span equation into
    Lemma~\ref{lem:block_restriction_intersection_subspace}, together with the
    normalization equation, gives the displayed intersection identity.
    Lemma~\ref{lem:product_span_block_image_direct_sum} at lengths \(n+1\) and
    \(n+2\) gives the two directness statements.
\end{proof}

\begin{lemma}[Normalized BNT blocks give an eventual block intersection]
    \label{lem:bnt_direct_sum_unital_block_intersection}
    \lean{MPSTensor.pgvwc07_iSup_restriction_intersection_eventually_of_bnt_directSum_unital_c1}
    \leanok
    \uses{lem:bnt_direct_sum_unital_block_intersection_ge}
    Let \(A^1,\ldots,A^r\) be separated BNT blocks satisfying the irreducibility,
    left-canonicality, and normalized self-overlap hypotheses used in the
    direct-sum comparison. Suppose \(L_0>0\), and for every block \(j\) the map
    \[
        \Gamma_{L_0}^{A^j}:M_{D_j}(\mathbb C)\to
        (\mathbb C^d)^{\otimes L_0}
    \]
    is injective. Assume also the normalization
    \[
        \sum_a A^j_aA^{j\,\dagger}_a=I .
    \]
    Then there is a length \(N\) such that, for every \(n\ge N\), with
    \[
        S_n:=\bigvee_jG_{n+1}(A^j),
    \]
    one has
    \[
        \left\{\psi:\forall b,\ \psi(-,b)\in S_n\right\}
        \cap
        \left\{\psi:\forall a,\ \psi(a,-)\in S_n\right\}
        =
        \bigvee_jG_{n+2}(A^j).
    \]
\end{lemma}

\begin{proof}
    \leanok
    Apply Lemma~\ref{lem:bnt_direct_sum_unital_block_intersection_ge} and take
    \[
        N=(L_0+1)+(r-1)3(L_0+1).
    \]
\end{proof}

\begin{lemma}[Block-diagonal boundary-condition formula]
    \label{lem:block_diagonal_boundary_condition_sum}
    \lean{MPSTensor.BlockSumGroundSpace.groundSpaceMap_toTensorFromBlocks_eq_sum_blockDiagonal}
    \leanok
    \uses{def:ground_space_parent}
    Let
    \[
        B=\bigoplus_{j=1}^g\mu_jA_j.
    \]
    For block-diagonal boundary conditions \(X_j\),
    \[
        \Gamma_L^B\!\left(\bigoplus_{j=1}^g X_j\right)
        =
        \sum_{j=1}^g\Gamma_L^{A_j}(\mu_j^LX_j).
    \]
\end{lemma}

\begin{proof}
    \leanok
    This is the diagonal-block trace formula specialized to a block-diagonal
    boundary matrix.  The \(j\)-th diagonal block of
    \(\bigoplus_kX_k\) is \(X_j\), giving the displayed sum.
\end{proof}

\begin{lemma}[One block in a block-diagonal local ground space]
    \label{lem:local_ground_space_block_le_assembled}
    \lean{MPSTensor.groundSpace_block_le_toTensorFromBlocks}
    \leanok
    \uses{def:ground_space_parent}
    Let
    \[
        B=\bigoplus_{k=1}^g\mu_kA_k.
    \]
    If \(\mu_j\ne0\), then, for every length \(L\),
    \[
        G_L(A_j)\subseteq G_L(B).
    \]
\end{lemma}

\begin{proof}
    \leanok
    A vector in \(G_L(A_j)\) has the form \(\Gamma_L^{A_j}(X)\).  Put
    \(\mu_j^{-L}X\) in the \(j\)-th diagonal block of a boundary matrix for
    \(B\) and put zero in every other block.  Applying \(\Gamma_L^B\) gives
    back \(\Gamma_L^{A_j}(X)\).
\end{proof}

\begin{lemma}[Local ground space of a block-diagonal tensor]
    \label{lem:block_sum_local_ground_space}
    \lean{MPSTensor.groundSpace_toTensorFromBlocks_eq_iSup}
    \leanok
    \uses{def:ground_space_parent, lem:block_diagonal_boundary_condition_sum,
        lem:local_ground_space_block_le_assembled}
    Let
    \[
        B=\bigoplus_{j=1}^g\mu_jA_j,
        \qquad
        \mu_j\ne0\quad(1\le j\le g).
    \]
    Then, for every length \(L\),
    \[
        G_L(B)=\bigvee_{j=1}^gG_L(A_j).
    \]
\end{lemma}

\begin{proof}
    \leanok
    A boundary matrix \(X\) for \(B\) has diagonal blocks \(X_j\), and
    block-diagonality gives
    \[
        \Gamma_L^B(X)
        =
        \sum_{j=1}^g\Gamma_L^{A_j}(\mu_j^LX_j).
    \]
    This proves \(G_L(B)\subseteq\bigvee_jG_L(A_j)\). The reverse inclusion is
    Lemma~\ref{lem:local_ground_space_block_le_assembled} for each block \(j\).
    The displayed equality follows.
\end{proof}

\begin{lemma}[Periodic constraints for a block-diagonal tensor]
    \label{lem:block_diagonal_periodic_constraints_propagated_sum}
    \lean{MPSTensor.chainGroundSpace_toTensorFromBlocks_le_iSup_groundSpace}
    \leanok
    \uses{lem:block_sum_local_ground_space,
        lem:periodic_constraints_propagated_subspace_sequence}
    Let
    \[
        B=\bigoplus_{j=1}^g\mu_jA_j,
        \qquad \mu_j\ne0\quad(1\le j\le g).
    \]
    Set
    \[
        S_M:=\bigvee_{j=1}^gG_M(A_j).
    \]
    If, for every \(M\ge L\),
    \[
        \mathbb C^d\otimes S_M\cap S_M\otimes\mathbb C^d=S_{M+1},
    \]
    then
    \[
        \mathcal G_{N,L}(B)\subseteq S_N.
    \]
\end{lemma}

\begin{proof}
    \leanok
    Lemma~\ref{lem:block_sum_local_ground_space} gives \(G_L(B)=S_L\).
    Apply Lemma~\ref{lem:periodic_constraints_propagated_subspace_sequence}
    to the sequence \(S_M\).
\end{proof}

\begin{lemma}[Periodic constraints lie in the block local-space sum]
    \label{lem:bnt_block_diagonal_periodic_constraints_propagated_sum}
    \lean{MPSTensor.chainGroundSpace_toTensorFromBlocks_le_iSup_groundSpace_of_ge_of_bnt_directSum_unital_c1}
    \leanok
    \uses{lem:bnt_direct_sum_unital_block_intersection_ge,
        lem:block_diagonal_periodic_constraints_propagated_sum}
    Let \(A_1,\ldots,A_g\) be separated normalized BNT blocks. Suppose
    \(L_0>0\), and for every block \(j\) the map
    \[
        \Gamma_{L_0}^{A^j}:M_{D_j}(\mathbb C)\to
        (\mathbb C^d)^{\otimes L_0}
    \]
    is injective. Assume also the normalization
    \[
        \sum_a A^j_aA^{j\,\dagger}_a=I .
    \]
    For nonzero scalars \(\mu_j\), set
    \[
        B=\bigoplus_{j=1}^g\mu_jA_j,
        \qquad
        S_M:=\bigvee_{j=1}^gG_M(A_j).
    \]
    If
    \[
        L\ge (L_0+1)+(g-1)3(L_0+1)+1,
    \]
    then, for every \(N\ge L\),
    \[
        \mathcal G_{N,L}(B)\subseteq S_N.
    \]
\end{lemma}

\begin{proof}
    \leanok
    Put
    \[
        T=(L_0+1)+(g-1)3(L_0+1).
    \]
    For \(M\ge L\), one has \(M-1\ge T\). Applying
    Lemma~\ref{lem:bnt_direct_sum_unital_block_intersection_ge} at
    \(n=M-1\) gives
    \[
        \mathbb C^d\otimes S_M\cap S_M\otimes\mathbb C^d=S_{M+1}
    \]
    for every \(M\ge L\). Lemma~\ref{lem:block_diagonal_periodic_constraints_propagated_sum}
    then propagates the local constraint from \(L\) to \(N\).
\end{proof}

\begin{lemma}[Periodic constraints and internality of the block local-space sum]
    \label{lem:bnt_block_diagonal_periodic_constraints_internal_sum}
    \lean{MPSTensor.chainGroundSpace_toTensorFromBlocks_le_iSup_and_iSupIndep_of_bnt_unital_c1}
    \leanok
    \uses{lem:bnt_block_diagonal_periodic_constraints_propagated_sum,
        lem:product_span_block_image_direct_sum}
    With the hypotheses and notation of
    Lemma~\ref{lem:bnt_block_diagonal_periodic_constraints_propagated_sum}, set
    \[
        S_N:=\bigvee_{j=1}^gG_N(A_j).
    \]
    Then, for every $N\ge L$,
    \[
        \mathcal G_{N,L}(B)\subseteq S_N,
    \]
    and the sum defining \(S_N\) is internal: if
    \[
        \phi_j\in G_N(A_j),
        \qquad
        \sum_{j=1}^g\phi_j=0,
    \]
    then $\phi_j=0$ for every $j$.
\end{lemma}

\begin{proof}
    \leanok
    The inclusion is
    Lemma~\ref{lem:bnt_block_diagonal_periodic_constraints_propagated_sum}. Since
    $N\ge L$ and
    \[
        L\ge (L_0+1)+(g-1)3(L_0+1)+1,
    \]
    the product-span directness range of
    Lemma~\ref{lem:product_span_block_image_direct_sum} applies at length
    $N$:
    \[
        \operatorname{span}\{(A^1_w,\ldots,A^g_w): |w|=N\}
        =
        \prod_j M_{D_j}(\mathbb C)
        \quad\Longrightarrow\quad
        \ker\!\left(
            \bigoplus_{j=1}^g G_N(A_j)
            \longrightarrow (\mathbb C^d)^{\otimes N},
            \quad
            (\phi_j)_j\longmapsto \sum_j\phi_j
        \right)=0.
    \]
\end{proof}

\begin{lemma}[Periodic chain vectors decompose in the local block spaces]
    \label{lem:bnt_block_diagonal_open_boundary_decomposition}
    \lean{MPSTensor.exists_unique_sum_groundSpace_of_chainGroundSpace_toTensorFromBlocks_of_bnt_unital_c1}
    \leanok
    \uses{lem:bnt_block_diagonal_periodic_constraints_propagated_sum}
    With the hypotheses and notation of
    Lemma~\ref{lem:bnt_block_diagonal_periodic_constraints_propagated_sum}, every
    vector
    \[
        \psi\in
        \mathcal G_{N,L}\!\left(\bigoplus_{j=1}^g\mu_jA_j\right)
    \]
    has a unique family \((\psi_j)_{j=1}^g\) such that
    \[
        \psi=\sum_{j=1}^g\psi_j,
        \qquad
        \psi_j\in G_N(A_j).
    \]
\end{lemma}

\begin{proof}
    \leanok
    \uses{lem:bnt_block_diagonal_periodic_constraints_internal_sum}
    Lemma~\ref{lem:bnt_block_diagonal_periodic_constraints_internal_sum} gives
    \[
        \mathcal G_{N,L}\!\left(\bigoplus_{j=1}^g\mu_jA_j\right)
        \subseteq
        \bigvee_{j=1}^g G_N(A_j),
    \]
    and the spaces \(G_N(A_j)\) are independent: for every \(j\) and every
    \(x_j\in G_N(A_j)\),
    \[
        \sum_j x_j=0
        \quad\Longrightarrow\quad
        x_j=0.
    \]
    Thus every
    \(\psi\in\mathcal G_{N,L}(\bigoplus_j\mu_jA_j)\) has a unique
    decomposition \(\psi=\sum_j\psi_j\) with \(\psi_j\in G_N(A_j)\).
    The remaining boundary step in \cite[Theorem~2blocks.2]{PerezGarcia2007Matrix}
    is to prove
    \[
        \psi_j\in\mathcal G_{N,L}(A_j)
    \]
    for each \(j\), by closing the periodic boundary with block-diagonal
    boundary conditions.
\end{proof}

\begin{lemma}[Block-diagonal boundary conditions for periodic chain vectors]
    \label{lem:bnt_block_diagonal_open_boundary_matrices}
    \lean{MPSTensor.exists_blockDiagonal_boundary_of_chainGroundSpace_toTensorFromBlocks_of_bnt_unital_c1}
    \leanok
    \uses{lem:bnt_block_diagonal_periodic_constraints_propagated_sum}
    With the hypotheses and notation of
    Lemma~\ref{lem:bnt_block_diagonal_periodic_constraints_propagated_sum}, every
    vector
    \[
        \psi\in
        \mathcal G_{N,L}\!\left(\bigoplus_{j=1}^g\mu_jA_j\right)
    \]
    admits block-diagonal boundary conditions \(X_j\) such that
    \[
        \psi=\Gamma_N^B\!\left(\bigoplus_{j=1}^gX_j\right),
        \qquad
        \Gamma_N^{A_j}(\mu_j^NX_j)\in G_N(A_j)
    \]
    for every \(j\).
\end{lemma}

\begin{proof}
    \leanok
    \uses{lem:bnt_block_diagonal_open_boundary_decomposition,
        lem:block_diagonal_boundary_condition_sum}
    Lemma~\ref{lem:bnt_block_diagonal_open_boundary_decomposition} gives
    \[
        \psi=\sum_{j=1}^g\phi_j,
        \qquad
        \phi_j\in G_N(A_j).
    \]
    Choose \(Y_j\) with \(\phi_j=\Gamma_N^{A_j}(Y_j)\), and set
    \(X_j=\mu_j^{-N}Y_j\). Since \(\mu_j\ne0\),
    \[
        \Gamma_N^{A_j}(\mu_j^NX_j)=\phi_j.
    \]
    The block-diagonal boundary formula then gives
    \[
        \Gamma_N^B\!\left(\bigoplus_{j=1}^gX_j\right)
        =
        \sum_{j=1}^g \Gamma_N^{A_j}(\mu_j^NX_j)
        =
        \psi.
    \]
    This proves the displayed statement.
\end{proof}

The lemma gives membership in the open-boundary space \(G_N(A_j)\). The later
boundary-closing step concerns the separate assertion
\[
    \Gamma_N^{A_j}(\mu_j^N X_j)\in\mathcal G_{N,L}(A_j).
\]

\begin{lemma}[Cyclic restrictions of block-diagonal boundary vectors]
    \label{lem:block_diagonal_boundary_cyclic_restriction_sum}
    \lean{MPSTensor.cyclicRestrictₗ_groundSpaceMap_toTensorFromBlocks_blockDiagonal_eq_sum}
    \leanok
    \uses{lem:block_diagonal_boundary_condition_sum}
    Let \(B=\bigoplus_{j=1}^g\mu_jA_j\), and fix block-diagonal boundary
    conditions \(X_j\). For every cyclic restriction \(R_{i,\tau}\) from
    \(N\) sites to \(L\) sites,
    \[
        R_{i,\tau}\!\left(
            \Gamma_N^B\!\left(\bigoplus_{j=1}^gX_j\right)
        \right)
        =
        \sum_{j=1}^g
        R_{i,\tau}\!\left(\Gamma_N^{A_j}(\mu_j^NX_j)\right).
    \]
\end{lemma}

\begin{proof}
    \leanok
    \uses{lem:block_diagonal_boundary_condition_sum}
    Apply Lemma~\ref{lem:block_diagonal_boundary_condition_sum} at length \(N\),
    then use linearity of the cyclic restriction \(R_{i,\tau}\).
\end{proof}

\begin{lemma}[Local direct-sum constraint from block-diagonal boundary conditions]
    \label{lem:block_diagonal_boundary_local_sum_mem}
    \lean{MPSTensor.blockDiagonal_boundary_cyclicRestrict_sum_mem_iSup_groundSpace}
    \leanok
    \uses{lem:block_diagonal_boundary_cyclic_restriction_sum,
        lem:block_sum_local_ground_space}
    Let \(B=\bigoplus_{j=1}^g\mu_jA_j\), with \(\mu_j\ne0\), and suppose
    \[
        \psi=\Gamma_N^B\!\left(\bigoplus_{j=1}^gX_j\right),
        \qquad
        \psi\in\mathcal G_{N,L}(B).
    \]
    Then every local constraint of \(\psi\) gives
    \[
        \sum_{j=1}^g
        R_{i,\tau}\!\left(\Gamma_N^{A_j}(\mu_j^NX_j)\right)
        \in
        \bigvee_{j=1}^gG_L(A_j).
    \]
\end{lemma}

\begin{proof}
    \leanok
    \uses{lem:block_diagonal_boundary_cyclic_restriction_sum,
        lem:block_sum_local_ground_space}
    The hypothesis \(\psi\in\mathcal G_{N,L}(B)\) says that the left-hand side
    before decomposition lies in \(G_L(B)\).  Lemma
    \ref{lem:block_diagonal_boundary_cyclic_restriction_sum} rewrites this local
    vector as the displayed sum, and Lemma~\ref{lem:block_sum_local_ground_space}
    identifies \(G_L(B)\) with \(\bigvee_jG_L(A_j)\).

    The next source step is not this membership statement, but the blockwise
    extraction
    \[
        A^j_{i_{m+1}}C^j_{i_1}=D^j_{i_{m+1}}A^j_{i_1}
    \]
    from \cite[Theorem~2blocks.2]{PerezGarcia2007Matrix}.
\end{proof}

\begin{lemma}[Cyclic windows not crossing the cut]
    \label{lem:block_diagonal_boundary_nonwrapping_component}
    \lean{MPSTensor.blockDiagonal_boundary_cyclicRestrict_component_mem_groundSpace_of_nonwrapping}
    \leanok
    \uses{def:ground_space_parent, rem:cyclic_window_operators}
    Fix block-diagonal boundary conditions $X_j$. If the cyclic window beginning
    at $i$ stays inside the chosen linear interval, $i+L\le N$, then
    for every block $j$,
    \[
        R_{i,\tau}\!\left(\Gamma_N^{A_j}(\mu_j^NX_j)\right)
        \in G_L(A_j).
    \]
\end{lemma}

\begin{proof}
    \leanok
    \uses{def:ground_space_parent, rem:cyclic_window_operators}
    When the window stays inside the chosen linear interval, the boundary matrix
    remains outside the window. Let $A_{\rm left}$ be the product on the sites
    before the window and $A_{\rm right}$ the product on the sites after the
    window. Trace cyclicity
    gives
    \[
        R_{i,\tau}\!\left(\Gamma_N^{A_j}(\mu_j^NX_j)\right)
        =
        \Gamma_L^{A_j}\!\left(A_{\rm right}\,\mu_j^NX_j\,A_{\rm left}\right).
    \]
    Hence the restricted vector lies in $G_L(A_j)$. If the cyclic interval
    crosses the boundary cut, the matrix $\mu_j^N X_j$ lies between the two
    pieces of the interval after trace rotation. The missing block-diagonal
    boundary-closing comparison is the identity
    \[
        A^j_{i_{m+1}}C^j_{i_1}=D^j_{i_{m+1}}A^j_{i_1}
    \]
    from \cite[Theorem~2blocks.2]{PerezGarcia2007Matrix}.
\end{proof}

\begin{lemma}[Boundary-crossing cyclic intervals from the boundary matrix identity]
    \label{lem:block_diagonal_boundary_crossing_matrix_identity}
    \lean{MPSTensor.blockDiagonal_boundary_cyclicRestrict_component_mem_groundSpace_of_crossing_matrix}
    \leanok
    \uses{def:ground_space_parent, rem:cyclic_window_operators}
    Assume \(0<N\), \(L\le N\), and let the cyclic interval beginning at \(i\)
    cross the boundary cut, so \(N<i+L\). Put \(a=i+L-N\). Suppose that there is
    a matrix \(E\) such that for every word \(\beta\) of length \(a\),
    \[
        \mu_j^N X_j A^j_\beta
        A^j_{\tau_a}\cdots A^j_{\tau_{i-1}}
        =
        A^j_\beta E.
    \]
    Then
    \[
        R_{i,\tau}\!\left(\Gamma_N^{A_j}(\mu_j^NX_j)\right)
        \in G_L(A_j).
    \]
\end{lemma}

\begin{proof}
    \leanok
    \uses{def:ground_space_parent, rem:cyclic_window_operators}
    Write a length-\(L\) word in the boundary-crossing cyclic interval as
    \(\alpha\beta\), where \(\alpha\) is the segment inside
    \(\{i,\ldots,N-1\}\) and \(\beta\) is the segment after the interval wraps
    past the cut, inside \(\{0,\ldots,a-1\}\). The full
    \(N\)-site word appearing in the restriction is then
    \[
        \beta\,\tau_a\cdots\tau_{i-1}\,\alpha .
    \]
    Trace cyclicity gives
    \[
        \operatorname{tr}\!\left(
        A^j_\beta A^j_{\tau_a}\cdots A^j_{\tau_{i-1}}
        A^j_\alpha \mu_j^NX_j\right)
        =
        \operatorname{tr}\!\left(
        A^j_\alpha \mu_j^NX_j A^j_\beta
        A^j_{\tau_a}\cdots A^j_{\tau_{i-1}}\right).
    \]
    The assumed boundary matrix identity changes the last expression to
    \[
        \operatorname{tr}\!\left(A^j_\alpha A^j_\beta E\right),
    \]
    which is the coefficient of \(\Gamma_L^{A_j}(E)\) at the word
    \(\alpha\beta\). Hence the restricted vector lies in \(G_L(A_j)\).
\end{proof}

\begin{lemma}[Last crossing-window trace form]
    \label{lem:block_diagonal_last_crossing_trace_form}
    \lean{MPSTensor.blockDiagonal_boundary_last_cyclicRestrict_sum_eq_leftBoundaryComponents}
    \leanok
    \uses{def:ground_space_map, rem:cyclic_window_operators}
    Write \(N=M+1\) and let the local length be \(m+2\le M+1\). Consider the
    cyclic interval beginning at the last site \(M\). For an outside
    configuration \(\tau\), put
    \[
        \rho=\tau_{m+1}\cdots\tau_{M-1},
        \qquad
        C^j_a=A^j_\rho A^j_a(\mu_j^{M+1}X_j).
    \]
    Then the block sum of the last crossing-window restrictions has the
    left-boundary trace form
    \[
        \sum_j
        R_{M,\tau}\!\left(\Gamma_{M+1}^{A_j}(\mu_j^{M+1}X_j)\right)
        =
        \sum_j \Phi^j,
        \qquad
        \Phi^j(a,w,b)=\operatorname{tr}\!\left(A^j_b C^j_a A^j_w\right).
    \]
\end{lemma}

\begin{proof}
    \leanok
    \uses{def:ground_space_map, rem:cyclic_window_operators}
    The cyclic interval beginning at \(M\) reads the last site first and then
    wraps to the sites \(0,\ldots,m\). Thus, for the local word \(awb\), the
    full word in the trace is
    \[
        w\,b\,\rho\,a .
    \]
    Hence the \(j\)-th coefficient is
    \[
        \operatorname{tr}\!\left(A^j_wA^j_bA^j_\rho A^j_a
        (\mu_j^{M+1}X_j)\right).
    \]
    Trace cyclicity rotates this to
    \[
        \operatorname{tr}\!\left(
        A^j_bA^j_\rho A^j_a(\mu_j^{M+1}X_j)A^j_w\right),
    \]
    which is the displayed left-boundary trace form.
\end{proof}

\begin{lemma}[Last crossing-window coefficient comparison]
    \label{lem:block_diagonal_last_crossing_coefficients_from_sum_membership}
    \lean{MPSTensor.blockDiagonal_boundary_last_coefficients_of_sum_mem_iSup}
    \leanok
    \uses{lem:block_diagonal_last_crossing_trace_form,
        lem:blockwise_boundary_identities_from_left_trace_membership}
    In the setting of Lemma~\ref{lem:block_diagonal_last_crossing_trace_form},
    assume that the length-\(m\) simultaneous word tuples span the product
    algebra and that
    \[
        \sum_j
        R_{M,\tau}\!\left(\Gamma_{M+1}^{A_j}(\mu_j^{M+1}X_j)\right)
        \in \bigvee_j G_{m+2}(A^j).
    \]
    Then there are matrices \(E_j\) such that, for every block \(j\) and all
    boundary letters \(a,b\),
    \[
        A^j_b\bigl(A^j_\rho A^j_a(\mu_j^{M+1}X_j)\bigr)
        =
        A^j_bE_jA^j_a .
    \]
\end{lemma}

\begin{proof}
    \leanok
    \uses{lem:block_diagonal_last_crossing_trace_form,
        lem:blockwise_boundary_identities_from_left_trace_membership}
    Lemma~\ref{lem:block_diagonal_last_crossing_trace_form} writes the local
    block sum in the form
    \[
        \Phi(a,w,b)=
        \sum_j\operatorname{tr}\!\left(A^j_bC^j_aA^j_w\right),
        \qquad
        C^j_a=A^j_\rho A^j_a(\mu_j^{M+1}X_j).
    \]
    Applying
    Lemma~\ref{lem:blockwise_boundary_identities_from_left_trace_membership}
    gives the displayed blockwise equations.
\end{proof}

\begin{lemma}[Last crossing-window coefficients from a block-diagonal chain vector]
    \label{lem:block_diagonal_last_crossing_coefficients_from_chain_ground_space}
    \lean{MPSTensor.blockDiagonal_boundary_last_coefficients_of_blockDiagonal_chainGroundSpace}
    \leanok
    \uses{lem:block_diagonal_boundary_local_sum_mem,
        lem:block_diagonal_last_crossing_coefficients_from_sum_membership}
    Let \(B=\bigoplus_j\mu_jA^j\), and assume \(\mu_j\ne0\) for every block
    \(j\). Suppose
    \[
        \psi=\Gamma_{M+1}^{B}\!\left(\bigoplus_jX_j\right),
        \qquad
        \psi\in\mathcal G_{M+1,m+2}(B).
    \]
    Assume also that the length-\(m\) simultaneous word tuples span the product
    algebra. Then, for every outside configuration \(\tau\), with
    \(\rho=\tau_{m+1}\cdots\tau_{M-1}\), there are matrices \(E_j\) such that
    \[
        A^j_b\bigl(A^j_\rho A^j_a(\mu_j^{M+1}X_j)\bigr)
        =
        A^j_bE_jA^j_a
    \]
    for every block \(j\) and all boundary letters \(a,b\).
\end{lemma}

\begin{proof}
    \leanok
    \uses{lem:block_diagonal_boundary_local_sum_mem,
        lem:block_diagonal_last_crossing_coefficients_from_sum_membership}
    The local constraint of \(\psi\) at the cyclic interval beginning at \(M\)
    gives
    \[
        \sum_j
        R_{M,\tau}\!\left(\Gamma_{M+1}^{A_j}(\mu_j^{M+1}X_j)\right)
        \in \bigvee_j G_{m+2}(A^j)
    \]
    by Lemma~\ref{lem:block_diagonal_boundary_local_sum_mem}. Applying
    Lemma~\ref{lem:block_diagonal_last_crossing_coefficients_from_sum_membership}
    gives the displayed equations.
\end{proof}

\begin{lemma}[Boundary-crossing intervals from a right boundary-matrix identity]
    \label{lem:block_diagonal_boundary_crossing_right_boundary_matrix}
    \lean{MPSTensor.blockDiagonal_boundary_cyclicRestrict_component_mem_groundSpace_of_boundary_identity}
    \leanok
    \uses{lem:block_diagonal_boundary_crossing_matrix_identity}
    Assume \(0<N\), \(L\le N\), and let the cyclic interval beginning at \(i\)
    cross the boundary cut, so \(N<i+L\). Put \(a=i+L-N\). Suppose that there is
    a matrix \(Y\) such that for every word \(\beta\) of length \(a\),
    \[
        \mu_j^N X_j A^j_\beta=A^j_\beta Y .
    \]
    Then
    \[
        R_{i,\tau}\!\left(\Gamma_N^{A_j}(\mu_j^NX_j)\right)
        \in G_L(A_j).
    \]
\end{lemma}

\begin{proof}
    \leanok
    \uses{lem:block_diagonal_boundary_crossing_matrix_identity}
    Let
    \[
        E:=Y A^j_{\tau_a}\cdots A^j_{\tau_{i-1}} .
    \]
    Then, for every word \(\beta\) on the segment \(0,\ldots,a-1\),
    \[
        \mu_j^N X_j A^j_\beta
        A^j_{\tau_a}\cdots A^j_{\tau_{i-1}}
        =
        A^j_\beta E.
    \]
    Lemma~\ref{lem:block_diagonal_boundary_crossing_matrix_identity} gives the
    stated local constraint.
\end{proof}

\begin{lemma}[Boundary-crossing cyclic intervals suffice]
    \label{lem:block_diagonal_boundary_component_crossing_windows_suffice}
    \lean{MPSTensor.blockDiagonal_boundary_component_chainGroundSpace_of_crossing_windows}
    \leanok
    \uses{def:chain_ground_space,
        lem:block_diagonal_boundary_nonwrapping_component}
    Assume \(0<N\) and \(L\le N\). Fix block-diagonal boundary conditions
    \(X_j\). Suppose that, for every block \(j\), every cyclic interval
    beginning at \(i\) with \(N<i+L\), and every configuration \(\tau\),
    \[
        R_{i,\tau}\!\left(\Gamma_N^{A_j}(\mu_j^NX_j)\right)
        \in G_L(A_j).
    \]
    Then, for every block \(j\),
    \[
        \Gamma_N^{A_j}(\mu_j^NX_j)\in \mathcal G_{N,L}(A_j).
    \]
\end{lemma}

\begin{proof}
    \leanok
    \uses{def:chain_ground_space,
        lem:block_diagonal_boundary_nonwrapping_component}
    By definition, membership in \(\mathcal G_{N,L}(A_j)\) is the collection of
    all cyclic length-\(L\) local constraints. For an interval beginning at \(i\),
    either \(i+L\le N\) or \(N<i+L\). In the first case use
    Lemma~\ref{lem:block_diagonal_boundary_nonwrapping_component}. In the second
    case use the displayed hypothesis.
\end{proof}

\begin{lemma}[Right boundary-matrix identities give componentwise periodic constraints]
    \label{lem:block_diagonal_boundary_component_right_boundary_matrices}
    \lean{MPSTensor.blockDiagonal_boundary_component_chainGroundSpace_of_boundary_identities}
    \leanok
    \uses{lem:block_diagonal_boundary_crossing_right_boundary_matrix,
        lem:block_diagonal_boundary_component_crossing_windows_suffice}
    Assume \(0<N\) and \(L\le N\). Fix block-diagonal boundary conditions
    \(X_j\). Suppose that, for every block \(j\), every boundary-crossing cyclic
    interval beginning at \(i\), and every configuration \(\tau\), there is a
    matrix \(Y_{j,i,\tau}\) such that, for every word \(\beta\) of length
    \(a=i+L-N\),
    \[
        \mu_j^N X_j A^j_\beta
        =
        A^j_\beta Y_{j,i,\tau}
    \]
    Then,
    for every block \(j\),
    \[
        \Gamma_N^{A_j}(\mu_j^NX_j)\in \mathcal G_{N,L}(A_j).
    \]
\end{lemma}

\begin{proof}
    \leanok
    \uses{lem:block_diagonal_boundary_crossing_right_boundary_matrix,
        lem:block_diagonal_boundary_component_crossing_windows_suffice}
    Lemma~\ref{lem:block_diagonal_boundary_crossing_right_boundary_matrix}
    turns the displayed boundary identity into the local constraint for each
    boundary-crossing interval.
    Lemma~\ref{lem:block_diagonal_boundary_component_crossing_windows_suffice}
    then yields
    \[
        \Gamma_N^{A_j}(\mu_j^NX_j)\in \mathcal G_{N,L}(A_j),
    \]
    handling both boundary-crossing intervals \(N<i+L\) and non-crossing
    intervals \(i+L\le N\).
\end{proof}

\begin{lemma}[Spanning complementary identities give componentwise constraints]
    \label{lem:block_diagonal_boundary_component_complementary_word_identities}
    \lean{MPSTensor.blockDiagonal_boundary_component_chainGroundSpace_of_complementary_word_identities}
    \leanok
    \uses{thm:word_span_common_boundary_matrix,
        lem:block_diagonal_boundary_component_right_boundary_matrices}
    Assume \(0<N\), \(L\le N\), and that for every block \(j\), the
    length-\(N-L\) word products of \(A_j\) span the full matrix algebra. Suppose
    that, for every block \(j\), every boundary-crossing cyclic interval
    beginning at \(i\), and every complementary word
    \(\rho\) of length \(N-L\), there is a matrix \(E_{j,i,\rho}\) such
    that for every word \(\beta\) of length \(i+L-N\),
    \[
        \mu_j^N X_j A^j_\beta A^j_\rho
        =
        A^j_\beta E_{j,i,\rho}.
    \]
    Then, for every block \(j\),
    \[
        \Gamma_N^{A_j}(\mu_j^NX_j)\in\mathcal G_{N,L}(A_j).
    \]
\end{lemma}

\begin{proof}
    \leanok
    \uses{thm:word_span_common_boundary_matrix,
        lem:block_diagonal_boundary_component_right_boundary_matrices}
    For fixed \(j,i\), apply
    Theorem~\ref{thm:word_span_common_boundary_matrix} to the family
    \[
        Z_\beta=\mu_j^NX_jA^j_\beta,\qquad F_\beta=A^j_\beta.
    \]
    Since the length-\(N-L\) complementary-word products span the full matrix
    algebra, there is a single matrix \(Y_{j,i}\) such that, for every \(\beta\),
    \[
        \mu_j^NX_jA^j_\beta=A^j_\beta Y_{j,i}
    \]
    Lemma~\ref{lem:block_diagonal_boundary_component_right_boundary_matrices}
    then gives the componentwise periodic-chain constraints.
\end{proof}

\begin{lemma}[Complementary identities under block injectivity give componentwise constraints]
    \label{lem:block_diagonal_boundary_component_complementary_word_identities_injective}
    \lean{MPSTensor.blockDiagonal_boundary_component_chainGroundSpace_of_complementary_word_identities_of_injective}
    \leanok
    \uses{lem:block_diagonal_boundary_component_complementary_word_identities,
        thm:wordspan_propagation_unital}
    Assume \(0<N\), \(L\le N\), and \(L+L_0\le N\). Suppose each block \(A_j\)
    is \(L_0\)-block-injective and satisfies the normalization equation
    \[
        \sum_a A^j_aA^{j\dagger}_a=I .
    \]
    Suppose also that the boundary-crossing matrix identities of
    Lemma~\ref{lem:block_diagonal_boundary_component_complementary_word_identities}
    hold for every complementary word \(\rho\) of length \(N-L\). Then, for
    every block \(j\),
    \[
        \Gamma_N^{A_j}(\mu_j^NX_j)\in\mathcal G_{N,L}(A_j).
    \]
\end{lemma}

\begin{proof}
    \leanok
    \uses{lem:block_diagonal_boundary_component_complementary_word_identities,
        thm:wordspan_propagation_unital}
    Since \(L+L_0\le N\), the complementary length satisfies \(L_0\le N-L\).
    Theorem~\ref{thm:wordspan_propagation_unital} therefore gives, for every
    block \(j\),
    \[
        S_{N-L}(A_j)=\MN{D_j}
    \]
    Apply
    Lemma~\ref{lem:block_diagonal_boundary_component_complementary_word_identities}.
\end{proof}

\begin{lemma}[The block-diagonal chain space lies between two block sums]
    \label{lem:bnt_block_diagonal_periodic_chain_inclusions}
    \lean{MPSTensor.chainGroundSpace_toTensorFromBlocks_two_inclusions_and_iSupIndep_of_bnt_unital_c1}
    \leanok
    \uses{lem:isup_chain_ground_space_block_le_assembled,
        lem:bnt_block_diagonal_periodic_constraints_internal_sum}
    With the hypotheses and notation of
    Lemma~\ref{lem:bnt_block_diagonal_periodic_constraints_propagated_sum}, one has
    \[
        \bigvee_{j=1}^g\mathcal G_{N,L}(A_j)
        \subseteq
        \mathcal G_{N,L}\!\left(\bigoplus_{j=1}^g\mu_jA_j\right)
        \subseteq
        \bigvee_{j=1}^gG_N(A_j),
    \]
    and the sum defining the right-hand side is internal.
\end{lemma}

\begin{proof}
    \leanok
    The first inclusion is
    Lemma~\ref{lem:isup_chain_ground_space_block_le_assembled}. The second
    inclusion and internal directness are
    Lemma~\ref{lem:bnt_block_diagonal_periodic_constraints_internal_sum}. The
    claimed statement is their conjunction.
    % The remaining arXiv:quant-ph/0608197 boundary-closing step with
    % block-diagonal boundary
    % conditions replaces
    % $\bigvee_jG_N(A_j)$ by $\sum_j\mathcal G_{N,L}(A_j)$.
\end{proof}

\begin{lemma}[Boundary inclusion gives the block-diagonal chain equality]
    \label{lem:block_diagonal_boundary_closing_equality}
    \lean{MPSTensor.chainGroundSpace_toTensorFromBlocks_eq_iSup_chainGroundSpace_of_boundary_closing}
    \leanok
    \uses{lem:isup_chain_ground_space_block_le_assembled}
    Let
    \[
        B=\bigoplus_{j=1}^g\mu_jA_j,
        \qquad \mu_j\ne0\quad(1\le j\le g).
    \]
    If closing the periodic boundary with block-diagonal boundary conditions gives
    \[
        \mathcal G_{N,L}(B)
        \subseteq
        \bigvee_{j=1}^g\mathcal G_{N,L}(A_j),
    \]
    then
    \[
        \mathcal G_{N,L}(B)
        =
        \bigvee_{j=1}^g\mathcal G_{N,L}(A_j).
    \]
\end{lemma}

\begin{proof}
    \leanok
    The reverse inclusion is
    Lemma~\ref{lem:isup_chain_ground_space_block_le_assembled}; combine the two
    inclusions.
\end{proof}

\begin{lemma}[Componentwise periodic decomposition gives the reverse inclusion]
    \label{lem:block_diagonal_boundary_decomposition_inclusion}
    \lean{MPSTensor.chainGroundSpace_toTensorFromBlocks_le_iSup_of_boundary_decomposition}
    \leanok
    Let
    \[
        B=\bigoplus_{j=1}^g\mu_jA_j.
    \]
    Suppose that every vector $\psi\in\mathcal G_{N,L}(B)$ can be written as
    \[
        \psi=\sum_{j=1}^g\psi_j,
        \qquad
        \psi_j\in\mathcal G_{N,L}(A_j).
    \]
    Then
    \[
        \mathcal G_{N,L}(B)
        \subseteq
        \bigvee_{j=1}^g\mathcal G_{N,L}(A_j).
    \]
\end{lemma}

\begin{proof}
    \leanok
    For $\psi\in\mathcal G_{N,L}(B)$, choose the displayed decomposition. Each
    summand lies in the corresponding subspace
    $\mathcal G_{N,L}(A_j)$, hence their finite sum lies in
    $\bigvee_j\mathcal G_{N,L}(A_j)$.
\end{proof}

\begin{lemma}[Block-diagonal boundary vectors in the periodic block sum]
    \label{lem:block_diagonal_boundary_vector_mem_chain_sum}
    \lean{MPSTensor.groundSpaceMap_toTensorFromBlocks_blockDiagonal_mem_iSup_chainGroundSpace}
    \leanok
    \uses{lem:block_diagonal_boundary_condition_sum}
    Let $B=\bigoplus_{j=1}^g\mu_jA_j$. Fix block-diagonal boundary conditions
    $X_j$.
    Suppose that, for every $j$,
    \[
        \Gamma_N^{A_j}(\mu_j^N X_j)\in\mathcal G_{N,L}(A_j).
    \]
    Then
    \[
        \Gamma_N^B\!\left(\bigoplus_{j=1}^gX_j\right)
        \in
        \bigvee_{j=1}^g\mathcal G_{N,L}(A_j).
    \]
\end{lemma}

\begin{proof}
    \leanok
    The block-diagonal boundary formula gives
    \[
        \Gamma_N^B\!\left(\bigoplus_{j=1}^gX_j\right)
        =
        \sum_{j=1}^g\Gamma_N^{A_j}(\mu_j^N X_j).
    \]
    Each summand belongs to the corresponding periodic block ground space by
    assumption, so the finite sum lies in the linear span of these spaces.
\end{proof}

\begin{lemma}[Block-diagonal boundary closing gives the reverse inclusion]
    \label{lem:block_diagonal_boundary_representation_inclusion}
    \lean{MPSTensor.chainGroundSpace_toTensorFromBlocks_le_iSup_of_blockDiagonal_boundary_groundSpaceMap}
    \leanok
    \uses{lem:block_diagonal_boundary_vector_mem_chain_sum}
    Let $B=\bigoplus_{j=1}^g\mu_jA_j$. Suppose every
    $\psi\in\mathcal G_{N,L}(B)$ admits block-diagonal boundary conditions
    $X_j$ such that
    \[
        \psi=\Gamma_N^B\!\left(\bigoplus_{j=1}^gX_j\right),
    \]
    and, for every $j$,
    \[
        \Gamma_N^{A_j}(\mu_j^N X_j)\in\mathcal G_{N,L}(A_j).
    \]
    Then
    \[
        \mathcal G_{N,L}(B)
        \subseteq
        \bigvee_{j=1}^g\mathcal G_{N,L}(A_j).
    \]
\end{lemma}

\begin{proof}
    \leanok
    Apply Lemma~\ref{lem:block_diagonal_boundary_vector_mem_chain_sum} to the
    block-diagonal boundary representation of each
    $\psi\in\mathcal G_{N,L}(B)$.
\end{proof}

\begin{lemma}[Block-diagonal boundary closing gives equality in the finite range]
    \label{lem:bnt_block_diagonal_boundary_representation_equality}
    \lean{MPSTensor.chainGroundSpace_toTensorFromBlocks_eq_iSup_and_iSupIndep_of_bnt_c1_blockBoundary}
    \leanok
    \uses{lem:block_diagonal_boundary_representation_inclusion,
        lem:block_diagonal_boundary_closing_equality,
        lem:bnt_block_diagonal_periodic_constraints_internal_sum}
    With the hypotheses and notation of
    Lemma~\ref{lem:bnt_block_diagonal_periodic_constraints_propagated_sum}, suppose
    every vector $\psi\in\mathcal G_{N,L}(B)$ has block-diagonal boundary
    conditions $X_j$ such that
    \[
        \psi=\Gamma_N^B\!\left(\bigoplus_{j=1}^gX_j\right),
    \]
    and, for every $j$,
    \[
        \Gamma_N^{A_j}(\mu_j^N X_j)\in\mathcal G_{N,L}(A_j).
    \]
    Then
    \[
        \mathcal G_{N,L}(B)=
        \bigvee_{j=1}^g\mathcal G_{N,L}(A_j),
    \]
    and the sum defining $\bigvee_jG_N(A_j)$ is internal.
\end{lemma}

\begin{proof}
    \leanok
    The block-diagonal boundary representation gives
    \[
        \mathcal G_{N,L}(B)
        \subseteq
        \bigvee_{j=1}^g\mathcal G_{N,L}(A_j)
    \]
    by Lemma~\ref{lem:block_diagonal_boundary_representation_inclusion}. The
    equality follows from Lemma~\ref{lem:block_diagonal_boundary_closing_equality};
    the internality statement is
    Lemma~\ref{lem:bnt_block_diagonal_periodic_constraints_internal_sum}.
\end{proof}

\begin{lemma}[Boundary decomposition gives equality in the finite injectivity range]
    \label{lem:bnt_block_diagonal_boundary_decomposition_equality}
    \lean{MPSTensor.chainGroundSpace_toTensorFromBlocks_eq_iSup_and_iSupIndep_of_bnt_c1_boundary_decomposition}
    \leanok
    \uses{lem:block_diagonal_boundary_decomposition_inclusion,
        lem:isup_chain_ground_space_block_le_assembled,
        lem:bnt_block_diagonal_periodic_constraints_internal_sum}
    Let $A_1,\ldots,A_g$ be separated normalized blocks in the basis of normal
    tensors, and let
    \[
        B=\bigoplus_{j=1}^g\mu_jA_j,
        \qquad \mu_j\ne0.
    \]
    Assume each block is injective at length $L_0$, assume
    $\sum_aA^j_aA^{j\dagger}_a=I$ for every $j$, and suppose
    \[
        L\ge (L_0+1)+(g-1)3(L_0+1)+1,
        \qquad
        N\ge L.
    \]
    If every vector $\psi\in\mathcal G_{N,L}(B)$ can be written as
    \[
        \psi=\sum_{j=1}^g\psi_j,
        \qquad
        \psi_j\in\mathcal G_{N,L}(A_j),
    \]
    then
    \[
        \mathcal G_{N,L}(B)
        =
        \bigvee_{j=1}^g\mathcal G_{N,L}(A_j),
    \]
    and the sum $\bigvee_jG_N(A_j)$ is internal.
\end{lemma}

\begin{proof}
    \leanok
    Lemma~\ref{lem:block_diagonal_boundary_decomposition_inclusion} gives
    \[
        \mathcal G_{N,L}(B)
        \subseteq
        \bigvee_{j=1}^g\mathcal G_{N,L}(A_j).
    \]
    The blockwise inclusion gives the reverse containment, hence equality. The
    internality of $\bigvee_jG_N(A_j)$ is the directness conclusion of
    Lemma~\ref{lem:bnt_block_diagonal_periodic_constraints_internal_sum}.
\end{proof}
